% !TEX root = ../ApproxDA-TransUNet.tex
\section{Related Work}
\label{sec:related_work}

\subsection{CNN-based Medical Image Segmentation}

Convolutional encoder-decoder architectures established the foundational paradigm for medical image segmentation, with U-Net~\cite{ronneberger2015unet} introducing the skip-connection design that underpins most segmentation networks, and successors~\cite{zhou2018unetpp,oktay2018attention,diakogiannis2020resunet,ibtehaz2020multiresunet} progressively adding dense connections, attention gating, and residual learning for improved feature reuse and training stability. These models offer strong computational efficiency and favorable inductive biases for local texture patterns, but their limited receptive fields make it difficult to model relationships between distant anatomical structures.

\subsection{Transformer-based Medical Image Segmentation}

Transformer architectures have significantly advanced medical image segmentation by enabling global dependency modeling beyond the locality of convolutional neural networks. Representative methods such as TransUNet~\cite{chen2021transunet}, Swin-Unet~\cite{cao2022swinunet}, UNETR~\cite{hatamizadeh2022unetr}, UCTransNet~\cite{wang2022uctransnet}, TransNorm~\cite{azad2022transnorm}, and DAEformer~\cite{azad2023daeformer} combine convolutional feature extraction with transformer-based contextual modeling and report strong segmentation performance across a wide range of medical imaging tasks. These methods typically apply the same attention design across tasks, and how the spatial scope of attention affects accuracy on different tasks is rarely examined.

\subsection{Dual Attention and DA-TransUNet}

DANet~\cite{fu2019dual} introduced Position Attention Modules (PAM) and Channel Attention Modules (CAM), capturing spatial and channel correlations at $\mathcal{O}(N^2C)$ and $\mathcal{O}(C^2N)$ complexity respectively. DA-TransUNet~\cite{sun2024datransunet} adapted this design for medical image segmentation by inserting dual-attention blocks at multiple skip connections of a hybrid R50-ViT-B/16 encoder-decoder, improving accuracy while inheriting the quadratic cost of full-map attention. DBAANet~\cite{li2025dbaanet} further extends dual-branch attention by aggregating position and channel attention across multiple scales, demonstrating the continued appeal of dual-attention paradigms for medical image segmentation.

\subsection{Efficient Attention and Adaptive Routing}

Windowed local attention~\cite{liu2021swin}, low-rank projection~\cite{wang2020linformer}, and grouped channel interaction~\cite{rahman2024emcad} have become standard strategies for reducing the cost of global attention, with learnable gating widely adopted to fuse multi-branch outputs~\cite{azad2023daeformer,rahman2024emcad}. DAMamba~\cite{ding2025damamba} further shows that dual-attention routing transfers effectively to state-space models for lightweight breast ultrasound segmentation, confirming the versatility of multi-branch attention beyond standard transformer architectures. Existing efficient attention methods are mainly evaluated by their computational savings, and their effect on accuracy is usually assumed to be small. Whether this holds equally across segmentation tasks with different spatial characteristics has received little controlled study. This work addresses that question by varying one approximation parameter at a time within a fixed architecture and training protocol.